\section{Experimental Evaluation}
\label{sec:experiment}

\subsection{Experimental Setup}

\textbf{Model and Hardware}. We evaluate on Llama-3-8B (32 layers, $d_{\text{model}}\!=\!4096$, $n_{\text{heads}}\!=\!32$, $n_{\text{kv\_heads}}\!=\!8$, $d_{\text{head}}\!=\!128$) with context length 16,400 and 128 decode steps at batch size 64.
GPU configuration is B100 (989 TFLOPS FP16, 3.35 TB/s HBM bandwidth).
PIM configuration: 2048 banks across 8 HBM3E stacks, 2KB row buffer, PE width 16B/cycle at 2GHz, 32 GB/s per-bank DRAM-to-rowbuffer bandwidth.
Asymmetric quantization uses zero-points $z_Q\!=\!1.0$, $z_K\!=\!2.0$ with per-channel quantization for keys.

\textbf{Simulator}. We implement a cycle-accurate C++ simulator extending LLMSimulator~\citep{duplex} with configurable near-bank PIM units.
The simulator reports per-stage timing breakdown (Q@K, softmax, score@V) and per-bank PIM component times (rowbuffer fill vs. PE compute).
All PIM latencies are verified against analytical cycle counts (ratio 1.00$\times$).

\textbf{Configurations}. We evaluate six configurations:
\begin{enumerate}[leftmargin=*,nosep]
\item \textbf{FP16 GPU}: All operations on GPU, float16 KV cache.
\item \textbf{FP16 PIM}: All attention on PIM (for reference), float16 KV.
\item \textbf{2-bit GPU}: Asymmetric 2-bit KV, all on GPU.
\item \textbf{2-bit Hybrid} (\pimattn{}): Q@K on PIM, softmax+score@V on GPU.
\item \textbf{2-bit All-PIM}: Both Q@K and score@V on PIM (baseline).
\item \textbf{2-bit Dequant PIM}: Dequantize to FP16 in PIM before GEMV.
\end{enumerate}

\subsection{Per-Step Attention Latency}

Table~\ref{tab:main} presents the per-step attention latency breakdown across all configurations.
All values are per decode step with 16 sequences per data-parallel group.

% TODO Table 2: Main results table — same as experiment_comparison.md

\begin{table}[t]
\caption{Per-step attention latency breakdown ($\mu s$). All PIM experiments use 2048 banks.}
\label{tab:main}
\centering
\small
\begin{tabular}{lcccccc}
\toprule
\textbf{Config} & \textbf{Q@K} & \textbf{Softmax} & \textbf{Score@V} & \textbf{AttnGen} & \textbf{pim\_rb} & \textbf{pim\_pe} \\
\midrule
FP16 GPU   & 97 & 0.00 & -- & 95  & -- & -- \\
FP16 PIM   & 66 & 0.03 & 66 & 132 & 131 & 66 \\
2-bit GPU  & 48 & 0.00 & -- & 48  & -- & -- \\
2-bit Hyb. & 33 & 0.01 & 34 & 67  & 33  & 17 \\
2-bit All  & 33 & 0.01 & 33 & 66  & 66  & 33 \\
2-bit Deq. & 66 & 0.01 & 66 & 132 & 66  & 99 \\
\bottomrule
\end{tabular}
\end{table}

\textbf{Key findings}:
\begin{itemize}
\item \textbf{\pimattn{} (2-bit Hybrid) outperforms GPU}: Per-sequence Q@K latency is $2.1\mu s$ vs. $3.0\mu s$ for 2-bit GPU — a $1.4\times$ speedup. Against FP16 GPU ($6.0\mu s$/seq), the advantage is $2.9\times$.
\item \textbf{Dequantize-then-compute is inefficient}: 2-bit Dequant PIM incurs $3\times$ PE time increase (99 vs. 33 $\mu s$), doubling total Q@K latency.
This validates our approach of computing directly on 2-bit data.
\item \textbf{PIM row-buffer time dominates}: In 2-bit configurations, $pim\_rb$ accounts for 50-66\% of total attention latency, indicating the computation is memory-bandwidth-bound rather than compute-bound at 2048 banks.
\item \textbf{Score@V benefits from GPU}: Comparing Hybrid vs. All-PIM, GPU execution of score@V achieves similar latency ($34\mu s$) to PIM ($33\mu s$) despite data movement overhead, due to GPU's superior float16 throughput.
\end{itemize}

\subsection{PIM Cycle Verification}

We verify all reported PIM latencies against analytical cycle-count models.
For the FP16 PIM configuration at $N_{\text{banks}}\!=\!2048$:
\begin{itemize}
\item Per-GEMV data: $(MK + KN + MN) \cdot 2\text{B} = 4.26\text{MB}$
\item Per-bank: $4.26\text{MB} / 2048 = 2082\text{B}$, requiring 2 row-buffer fills
\item Row-buffer time: $2 \times 64\text{ns} = 128\text{ns}$ per GEMV
\item PE time: $2082\text{B} / 16\text{B/cycle} \times 0.5\text{ns} = 65.1\text{ns}$ per GEMV
\item Per-step (Q@K + score@V, $\times 4$ group, $\times 8$ heads, $\times 16$ seqs): RB $131.1\mu s$, PE $66.6\mu s$
\end{itemize}
Reported values: $pim\_rb = 131.1\mu s$, $pim\_pe = 66.5\mu s$.
\textbf{Ratio: 1.00$\times$ for both RB and PE}, confirming the simulator's fidelity.

\subsection{Asymmetric Quantization Overhead}

Figure~\ref{fig:asym} illustrates the impact of asymmetric quantization overhead on per-GEMV latency.
% TODO Figure 6: Bar chart comparing per-GEMV latency with/without asymmetric quant for FP16 and 2-bit

With asymmetric quantization enabled, the additional reduction terms ($\sum\hat{q}$ and $\sum\hat{k}$) add PE time proportional to $(M+N) \cdot K$.
For decode mode ($M=1, N=16528$), this matches the main GEMV data volume, resulting in $\sim\!1.8\times$ PE time increase.
In prefill mode with larger $M$, the relative overhead diminishes.

\subsection{End-to-End Throughput}

Table~\ref{tab:e2e} reports end-to-end latency for 128 decode steps across 32 layers.

% TODO Table 3: End-to-end results

\begin{table}[t]
\caption{End-to-end latency ($\mu s$) for 128 steps $\times$ 32 layers.}
\label{tab:e2e}
\centering
\small
\begin{tabular}{lcc}
\toprule
\textbf{Config} & \textbf{Total} & \textbf{Per Layer} \\
\midrule
FP16 GPU   & 5,623  & 176 \\
FP16 PIM   & 6,788  & 212 \\
2-bit GPU  & 2,769  & 87  \\
2-bit Hyb. & 3,460  & 108 \\
2-bit All  & 3,443  & 108 \\
2-bit Deq. & 5,555  & 174 \\
\bottomrule
\end{tabular}
\end{table}

The 2-bit GPU configuration achieves the lowest total latency (2,769$\mu s$) due to the combination of reduced memory traffic and GPU's high throughput.
However, \pimattn{} (2-bit Hybrid) at 3,460$\mu s$ is only $1.25\times$ slower than 2-bit GPU while offering the architectural advantage of keeping KV cache in PIM memory, potentially enabling larger batch sizes and longer contexts that exceed GPU memory capacity.
